\begin{tikzpicture}[
  scale=1.1,
  every node/.style={font=\scriptsize},
  >=Latex
]
% Binary, merger, and post-merger black hole
\fill[black] (1.00,2.45) circle (.15);
\fill[black] (1.55,2.45) circle (.15);
\draw[gray,-{Latex[length=2mm]},line width=.45pt]
  (1.28,2.74)
  arc[start angle=90,end angle=360,x radius=.47,y radius=.28];
\node[gray] at (1.28,2.04) {binary};

\draw[gray,-{Latex[length=2mm]}]
  (2.00,2.45)--(2.72,2.45);

\fill[black] (3.16,2.45) circle (.23);
\fill[black] (3.46,2.45) circle (.23);
\draw[red,line width=.55pt]
  (3.00,2.45)
  .. controls (3.18,2.83) and (3.50,2.83) ..
  (3.66,2.45);
\node[red] at (3.31,2.04) {merger};

\draw[gray,-{Latex[length=2mm]}]
  (3.86,2.45)--(4.58,2.45);

\fill[black] (5.02,2.45) circle (.25);
\draw[blue,line width=.65pt]
  (4.64,2.45)
  arc[start angle=180,end angle=0,x radius=.38,y radius=.19];
\draw[blue!65,line width=.45pt]
  (4.56,2.45)
  arc[start angle=180,end angle=0,x radius=.46,y radius=.28];
\node[blue] at (5.02,2.04) {post-merger};

% Domains of validity
\fill[red!7]
  (3.45,-.92) rectangle (4.45,1.12);
\fill[blue!5]
  (4.45,-.92) rectangle (11.15,1.12);

\draw[red,line width=1.4pt]
  (3.45,1.31)--(4.45,1.31);
\node[red,above] at (3.95,1.31)
  {nonlinear general relativity};

\draw[blue,line width=1.4pt]
  (4.45,1.31)--(11.15,1.31);
\node[blue,above] at (7.80,1.31)
  {linear black-hole perturbation theory};
\node[blue,below,align=center] at (7.80,1.31)
  {$g_{\mu\nu}=g^{\mathrm{BH}}_{\mu\nu}+\delta g_{\mu\nu}$,
   $|\delta g|\ll|g^{\mathrm{BH}}|$};

% Waveform axes
\draw[gray,-{Latex[length=2mm]}]
  (.18,0)--(11.42,0)
  node[right] {$t$};
\draw[gray,-{Latex[length=2mm]}]
  (.34,-.92)--(.34,1.02)
  node[above] {$h(t)$};

\def\tpeak{3.48}
\def\tlin{4.45}
\def\ttail{8.10}

% Match amplitude and phase at the joining points.
\pgfmathsetmacro{\Apeak}{.055*exp(.60*\tpeak)}
\pgfmathsetmacro{\phipeak}{74*\tpeak+25*\tpeak*\tpeak}
\pgfmathsetmacro{\Alin}{\Apeak*exp(-.10*(\tlin-\tpeak))}
\pgfmathsetmacro{\philin}{\phipeak+395*(\tlin-\tpeak)}
\pgfmathsetmacro{\Atail}{
  \Alin
  *exp(-.50*(\ttail-\tlin))
  *sin(\philin+420*(\ttail-\tlin))
}

% Inspiral
\draw[
  black,line width=.75pt,smooth,
  domain=.45:\tpeak,samples=220
]
  plot
  (\x,{
    .055*exp(.60*\x)
    *sin(74*\x+25*\x*\x)
  });

% Nonlinear merger
\draw[
  black,line width=.95pt,smooth,
  domain=\tpeak:\tlin,samples=100
]
  plot
  (\x,{
    \Apeak
    *exp(-.10*(\x-\tpeak))
    *sin(\phipeak+395*(\x-\tpeak))
  });

% Exponential envelope of the QNM ringdown
\draw[
  red!60,dashed,line width=.45pt,smooth,
  domain=\tlin:\ttail,samples=80
]
  plot
  (\x,{
    \Alin*exp(-.50*(\x-\tlin))
  });
\draw[
  red!60,dashed,line width=.45pt,smooth,
  domain=\tlin:\ttail,samples=80
]
  plot
  (\x,{
    -\Alin*exp(-.50*(\x-\tlin))
  });

% QNM-dominated ringdown
\draw[
  red,line width=1.05pt,smooth,
  domain=\tlin:\ttail,samples=220
]
  plot
  (\x,{
    \Alin
    *exp(-.50*(\x-\tlin))
    *sin(\philin+420*(\x-\tlin))
  });

% Late-time power-law tail
\draw[
  blue,line width=1.05pt,smooth,
  domain=\ttail:11.15,samples=100
]
  plot
  (\x,{
    \Atail/pow(1+.62*(\x-\ttail),2.4)
  });

% Transition times
\draw[gray,densely dotted]
  (3.48,-1.02)--(3.48,1.05);
\draw[blue,densely dashed]
  (4.45,-1.02)--(4.45,1.12);
\draw[gray,densely dotted]
  (8.10,-1.02)--(8.10,1.05);

\node[below,gray] at (3.48,-1.00)
  {$t_{\mathrm{peak}}$};
\node[below,blue] at (4.45,-1.00)
  {$t_{\mathrm{lin}}$};
\node[below,gray] at (8.10,-1.00)
  {$t_{\mathrm{tail}}$};

% Ringdown and tail
\draw[red,line width=1.0pt]
  (4.58,-1.46)--(7.96,-1.46);
\node[below,red] at (6.24,-1.46)
  {QNM-dominated ringdown};

\draw[blue,line width=1.0pt]
  (8.24,-1.46)--(11.03,-1.46);
\node[below,blue,align=center] at (9.64,-1.46)
  {tail: $h\sim t^{-p}$, $p>0$\\
   branch cut / continuous component};
\end{tikzpicture}
